% Generated draft: requires mathematical and source review.
\documentclass[11pt]{article}
\usepackage[margin=25mm]{geometry}
\usepackage{fontspec}
\setmainfont{Latin Modern Roman}
\usepackage{amsmath,amssymb,mathtools,unicode-math}
\setmathfont{Latin Modern Math}
\usepackage{enumitem,needspace,xurl,hyperref}
\setlength{\parindent}{0pt}
\setlength{\parskip}{5pt}
\setlength{\emergencystretch}{3em}
\newcommand{\N}{\mathbb{N}}
\newcommand{\Z}{\mathbb{Z}}
\newcommand{\Q}{\mathbb{Q}}
\newcommand{\R}{\mathbb{R}}
\newcommand{\C}{\mathbb{C}}
\newcommand{\F}{\mathbb{F}}
\newcommand{\A}{\mathbb{A}}
\newcommand{\PP}{\mathbb{P}}
\newcommand{\E}{\mathbb{E}}
\newcommand{\eps}{\varepsilon}
\newcommand{\dd}{\,\mathrm{d}}
\begin{document}
\section*{UnsolvedMath ID: 20001882}
\textbf{Source problem code:} catalogue entry
\subsection{Definitions and mathematical statement}
Let \(M\) be a closed connected Riemannian manifold without conjugate points, meaning that no nonzero Jacobi field along a geodesic vanishes at two distinct times. Let \(X=\widetilde M\) with its lifted complete metric and \(\Gamma=\pi_1(M)\). The source asks whether \(\Gamma\) is semihyperbolic in the sense of Alonso--Bridson. In a standard group formulation, this means that a Cayley graph admits an equivariant bicombing by uniformly quasigeodesic paths with a uniform synchronous endpoint-stability bound; equivalently the bound can be written affine in endpoint displacement. This is a large-scale group property, not just a pointwise uniqueness assertion for geodesics.

On \(X\), let \(b_{xy}:[0,\infty)\to X\) be the unique geodesic from \(x\) to \(y\), traversed at unit speed until time \(d(x,y)\), and constantly equal to \(y\) thereafter. The unique-geodesic construction is \(\Gamma\)-equivariant. Define
\[
 W(R)=\sup_{\substack{d(x,z)\leq R\\d(y,w)\leq R}}
       \sup_{t\geq0}d(b_{xy}(t),b_{zw}(t)).
\]
The aid's criterion states that the following are equivalent:
\[
 \exists r_0>0,\ C<\infty:\quad W(r_0)\leq C;
\]
\[
 \exists A,B<\infty:\quad
 \sup_{t\geq0}d(b_{xy}(t),b_{zw}(t))
 \leq A\max\{d(x,z),d(y,w)\}+B
 \quad\text{for all }x,y,z,w\in X.
\]
More explicitly, the first condition gives
\[
 W(R)\leq C\lceil R/r_0\rceil\leq (C/r_0)R+C
 \qquad(R>0)
\]
by subdividing both endpoint paths; \(W(0)=0\). Such a bound, together with the proper cocompact deck action and the geodesic bicombing, is sufficient for semihyperbolicity of \(\Gamma\). The aid does not establish that every no-conjugate-points metric has this canonical endpoint bound, nor claim that semihyperbolicity of the abstract group requires this particular bicombing to work.

The reported stronger-hypothesis result assumes monotone equal-radius divergence: whenever two unit-speed geodesic segments or rays have a common initial point, their mutual distance at equal parameters is nondecreasing as long as both parameters are defined. Under that property,
\[
 d(y,w)\leq R\quad\Longrightarrow\quad
 \sup_t d(b_{xy}(t),b_{xw}(t))\leq2R,
\]
\[
 d(x,z)\leq R\quad\Longrightarrow\quad
 \sup_t d(b_{xy}(t),b_{zy}(t))\leq3R.
\]
The extra term for a common terminal point accounts for the possibly unequal path lengths under unit-speed terminally constant parametrization. Consequently
\[
 d(x,z),d(y,w)\leq R\quad\Longrightarrow\quad
 \sup_{t\geq0}d(b_{xy}(t),b_{zw}(t))\leq5R.
\]
The aid states that manifolds without focal points have the required monotone divergence and hence satisfy this criterion. This gives a sufficient subclass and an explicit constant, while the general no-conjugate-points question remains open in the supplied account. Changing to a normalized \([0,1]\) parametrization would be a different synchronous estimate and is not being done here.
\subsection{Short English statement}
The question asks whether every closed manifold without conjugate points has a semihyperbolic fundamental group. The aid reduces this to a sufficient uniform endpoint bound for the canonical geodesic paths and proves an explicit \(5R\) bound under the stronger no-focal-points divergence property. It does not settle the general case.
\subsection{Sources}
\begin{itemize}
\item Permanent UnsolvedMath ID \(20001882\), exact public code catalogue entry; local registry \texttt{aim-geometry-notes.json}, \texttt{notes[219]}, canonical tag question.
\item AIM workshop ``Geodesics,'' Geometry, extracted source item 4.1, Question 4.1.3: \url{https://aimath.org/pastworkshops/geodesicsproblems.pdf}. No PDF was accessed. No proposer or year is supplied.
\item The original background attributes semihyperbolicity to Alonso and Bridson [Alo-Bri1995]. No full title or publication details are given, and the original literature field is otherwise empty.
\item Research attempt 1, ``Canonical endpoint stability and an explicit no-focal-points bound,'' labelled partial\_result, open at its run, proved and low novelty confidence; registry partially\_solved. The exact-statement verification field identifies the canonical question but supplies no independent external verification evidence.
\end{itemize}
\end{document}
